\documentclass[10pt, aspectratio=169]{beamer}

% ==========================================
% Theme & Color Settings (Moloch Theme)
% ==========================================
\usetheme{moloch}
\usecolortheme{moloch}

% ==========================================
% Modular Preambles & Configurations (Relative)
% ==========================================
\input{../../preamble/packages.tex}
\input{../../preamble/macros.tex}
\input{../../preamble/listings-setup.tex}

% ==========================================
% Presentation Metadata
% ==========================================
\title{Algorithmic Codesign of Federated Distillation}
\subtitle{Research Progress Update - Year 3, Trimester 3, Week 3}
\author{Christian Joseph Bunyi}
\institute{Dr. Andrew L. Tan Data Science Institute \\ De La Salle University}
\date{June 1, 2026}

% ==========================================
% Slide Deck Document Structure
% ==========================================
\begin{document}

% Slide 1: Title Slide
\begin{frame}[plain]
    \titlepage
\end{frame}

% Slide 2: Research Domain Convergence (TikZ 3-Pillar Diagram)
\begin{frame}{Research Domain Convergence}
    \centering
    \begin{tikzpicture}[node distance=1.2cm, auto, >=Stealth, thick]
    % Define theme-harmonious soft colors
    \definecolor{colorfd}{RGB}{111,168,220}      % Soft blue
    \definecolor{coloriot}{RGB}{147,196,125}     % Soft green
    \definecolor{coloropt}{RGB}{224,102,102}     % Soft red/rose
    \definecolor{colortarget}{RGB}{246,178,107}  % Soft orange

    % Three Pillar Nodes (Stacked/relative horizontal flow)
    \node[draw=colorfd, fill=colorfd!15, rounded corners, minimum width=3.2cm, minimum height=1.2cm, align=center] (fd) {
        \textbf{Federated Distillation}\\
        \fontsize{7.5pt}{9.5pt}\selectfont Gradient-Free Soft-Labels
    };

    \node[draw=coloriot, fill=coloriot!15, rounded corners, minimum width=3.2cm, minimum height=1.2cm, align=center, right=0.6cm of fd] (iot) {
        \textbf{Healthcare IoT (IoMT)}\\
        \fontsize{7.5pt}{9.5pt}\selectfont Device Resource Bounds
    };

    \node[draw=coloropt, fill=coloropt!15, rounded corners, minimum width=3.2cm, minimum height=1.2cm, align=center, right=0.6cm of iot] (opt) {
        \textbf{Adaptive Optimization}\\
        \fontsize{7.5pt}{9.5pt}\selectfont Client Data Heterogeneity
    };

    % Target Unified Focus Node (Proposed Co-design) centered below
    \node[draw=colortarget, fill=colortarget!15, rounded corners, minimum width=7.5cm, minimum height=1.3cm, align=center, below=1.2cm of iot] (target) {
        \textbf{Proposed Algorithmic Co-design}\\
        \fontsize{8.5pt}{10.5pt}\selectfont \keyterm{Hybrid Data- and Resource-Aware Quantization}
    };

    % Connecting Arrows
    \draw[->, thick, colorfd] (fd.south) -- (target.north west);
    \draw[->, thick, coloriot] (iot.south) -- (target.north);
    \draw[->, thick, coloropt] (opt.south) -- (target.north east);
\end{tikzpicture}

    \vspace{0.4cm}
    \begin{center}
        \fontsize{9pt}{11pt}\selectfont
        Convergence of soft-label communication, local hardware constraints, and data distribution complexity.
    \end{center}
\end{frame}

% Slide 3: Literature Landscape & Key Research Gaps
\begin{frame}{Literature Landscape \& Key Research Gaps}
    \begin{columns}[T]
        \column{0.45\textwidth}
        \textbf{The Compression-Utility Trade-off}
        \begin{itemize}
            \item \keyterm{Core Limitation:} Static compression baseline methods destroy fine-grained activity recognition features.
            \item \keyterm{Literature Gap:} Hardware-only dynamic triggers blindly over-compress critical motion data in skewed settings.
            \item \keyterm{Objective:} Balancing communication savings with patient fall-recognition accuracy.
        \end{itemize}

        \column{0.55\textwidth}
        \centering
        \textbf{Taxonomy Mind Map}
        \vspace{0.2cm}

        \begin{center}
            \setlength{\fboxrule}{0.5pt}
            \fbox{
                \begin{minipage}{0.9\textwidth}
                    \centering
                    \vspace{0.6cm}
                    \fontsize{7.5pt}{9.5pt}\selectfont
                    \texttt{[Placeholder: Miro Mind Map Graphic]} \\
                    \vspace{0.2cm}
                    \fontsize{6.5pt}{8.5pt}\selectfont
                    \textit{Source: Export full Miro canvas mind map detailing research dimensions}
                    \vspace{0.6cm}
                \end{minipage}
            }
        \end{center}
    \end{columns}
\end{frame}

% Slide 4: Domain 1 — Intelligent Edge Computing & IoT Optimization
\begin{frame}{Intelligent Edge \& IoT Optimization}
    \begin{columns}[T]
        \column{0.45\textwidth}
        \textbf{Domain Summary}
        \begin{itemize}
            \item \keyterm{Constraints:} High peak memory overhead during local training.
            \item \keyterm{Hardware Limits:} Strict energy and bandwidth exhaustions on wearables.
            \item \keyterm{Key Finding:} Fixed scheduling is highly suboptimal for variable edge networks.
        \end{itemize}

        \column{0.55\textwidth}
        \centering
        \textbf{Edge Architecture Hierarchy}
        \vspace{0.2cm}

        \begin{center}
            \setlength{\fboxrule}{0.5pt}
            \fbox{
                \begin{minipage}{0.9\textwidth}
                    \centering
                    \vspace{0.6cm}
                    \fontsize{7.5pt}{9.5pt}\selectfont
                    \texttt{[Placeholder: Intelligent Edge Systems]} \\
                    \vspace{0.2cm}
                    \fontsize{6.5pt}{8.5pt}\selectfont
                    \textit{Reference: Retrieve Figure 2 (Edge Scheduling Layers) from Cajas Ordonez et al. 2025}
                    \vspace{0.6cm}
                \end{minipage}
            }
        \end{center}
    \end{columns}
\end{frame}

% Slide 5: Domain 2 — Model Compression in Federated Learning
\begin{frame}{Model Compression in Federated Learning}
    \begin{columns}[T]
        \column{0.45\textwidth}
        \textbf{Compression Methods}
        \begin{itemize}
            \item \keyterm{Primary Tools:} Quantization, sparsification, and pruning algorithms.
            \item \keyterm{Gap:} Standard static compression destroys fine-grained local motion features.
            \item \keyterm{Need:} Dynamic feedback loops reflecting client data variance.
        \end{itemize}

        \column{0.55\textwidth}
        \centering
        \textbf{Compression Classification}
        \vspace{0.2cm}

        \begin{center}
            \setlength{\fboxrule}{0.5pt}
            \fbox{
                \begin{minipage}{0.9\textwidth}
                    \centering
                    \vspace{0.6cm}
                    \fontsize{7.5pt}{9.5pt}\selectfont
                    \texttt{[Placeholder: Compression Taxonomy]} \\
                    \vspace{0.2cm}
                    \fontsize{6.5pt}{8.5pt}\selectfont
                    \textit{Reference: Retrieve Figure 4 (Taxonomy of Compression in FL) from Le et al. 2024}
                    \vspace{0.6cm}
                \end{minipage}
            }
        \end{center}
    \end{columns}
\end{frame}

% Slide 6: Domain 3 — Federated Distillation (FD) Frameworks
\begin{frame}{Federated Distillation (FD) Frameworks}
    \begin{columns}[T]
        \column{0.45\textwidth}
        \textbf{Distillation Mechanics}
        \begin{itemize}
            \item \keyterm{Exchange:} Transfers soft-label logits rather than raw network weights.
            \item \keyterm{Bandwidth:} Drastically reduces transmission sizes on channels.
            \item \keyterm{Robustness:} Mitigates client statistical skew (non-IID) during aggregation.
        \end{itemize}

        \column{0.55\textwidth}
        \centering
        \textbf{Knowledge Distillation Flow}
        \vspace{0.2cm}

        \begin{center}
            \setlength{\fboxrule}{0.5pt}
            \fbox{
                \begin{minipage}{0.9\textwidth}
                    \centering
                    \vspace{0.6cm}
                    \fontsize{7.5pt}{9.5pt}\selectfont
                    \texttt{[Placeholder: Logit Distillation Flow]} \\
                    \vspace{0.2cm}
                    \fontsize{6.5pt}{8.5pt}\selectfont
                    \textit{Reference: Retrieve Figure 1 (Logit Aggregation Scheme) from Salman et al. 2025}
                    \vspace{0.6cm}
                \end{minipage}
            }
        \end{center}
    \end{columns}
\end{frame}

% Slide 7: AWS Certified Solutions Architect - Associate Certification
\begin{frame}{AWS Solutions Architect - Associate}
    \centering
    \vfill
    \begin{center}
        \setlength{\fboxrule}{0.5pt}
        \fbox{
            \begin{minipage}{0.6\textwidth}
                \centering
                \vspace{1.5cm}
                \fontsize{10pt}{12pt}\selectfont
                \texttt{[Placeholder: Official AWS Badge Image]} \\
                \vspace{0.4cm}
                \fontsize{8pt}{10pt}\selectfont
                \textit{Source: Official Credly AWS Certification Badge from the internet}
                \vspace{1.5cm}
            \end{minipage}
        }
    \end{center}
    \vfill
\end{frame}

% Slide 8: Workspace & Repository Organization (Split A)
\begin{frame}{Workspace \& Repository Organization}
    \begin{columns}[T]
        \column{0.45\textwidth}
        \textbf{Workspace Management}
        \begin{itemize}
            \item \keyterm{Collaboration:} Set up a dedicated, clean GitHub organization containing all research-level repositories.
            \item \keyterm{Consistency:} Aligns all environment scripts and code baselines across human and agentic developers.
            \item \keyterm{Reproducibility:} Establishes standardized folder conventions to support automated checks.
        \end{itemize}

        \column{0.55\textwidth}
        \centering
        \textbf{GitHub Repository Layout}
        \vspace{0.2cm}

        \begin{center}
            \setlength{\fboxrule}{0.5pt}
            \fbox{
                \begin{minipage}{0.9\textwidth}
                    \centering
                    \vspace{0.6cm}
                    \fontsize{7.5pt}{9.5pt}\selectfont
                    \texttt{[Placeholder: GitHub Organization Structure]} \\
                    \vspace{0.2cm}
                    \fontsize{6.5pt}{8.5pt}\selectfont
                    \textit{Source: Take screenshot of your GitHub organization listing repositories}
                    \vspace{0.6cm}
                \end{minipage}
            }
        \end{center}
    \end{columns}
\end{frame}

% Slide 9: Simulation Framework Selection - Flower (Split B)
\begin{frame}{Simulation Framework Selection}
    \begin{columns}[T]
        \column{0.45\textwidth}
        \textbf{Flower FL Framework}
        \begin{itemize}
            \item \keyterm{Modularity:} Selected \texttt{Flower} due to its clean client-side abstractions.
            \item \keyterm{IoT Scaling:} Lightweight local orchestrations, perfectly reflecting wearable memory limitations.
            \item \keyterm{Heterogeneity:} Seamlessly simulates statistical skew on client nodes.
        \end{itemize}

        \column{0.55\textwidth}
        \centering
        \textbf{Flower Architecture}
        \vspace{0.2cm}

        \begin{center}
            \setlength{\fboxrule}{0.5pt}
            \fbox{
                \begin{minipage}{0.9\textwidth}
                    \centering
                    \vspace{0.6cm}
                    \fontsize{7.5pt}{9.5pt}\selectfont
                    \texttt{[Placeholder: Flower Core Architecture Diagram]} \\
                    \vspace{0.2cm}
                    \fontsize{6.5pt}{8.5pt}\selectfont
                    \textit{Source: Grab screenshot of Flower architecture diagram from flower.ai}
                    \vspace{0.6cm}
                \end{minipage}
            }
        \end{center}
    \end{columns}
\end{frame}

% Slide 10: Baseline Replication & Validation (Split C)
\begin{frame}[fragile]{Baseline Replication \& Validation}
    \begin{columns}[T]
        \column{0.48\textwidth}
        \textbf{Local Baseline Execution}
        \begin{itemize}
            \item \keyterm{FedAvg \& FedProx:} Replicated both classic baseline models successfully on local hardware.
            \item \keyterm{Parameter Tuning:} Configured SGD optimizers and proximal penalty metrics.
        \end{itemize}

        \begin{table}
            \centering
            \fontsize{6.5pt}{8pt}\selectfont
            \caption{Local Baseline Replication Setup}
            \begin{tabular}{lcc}
                \toprule
                \textbf{Parameter}       & \textbf{FedAvg} & \textbf{FedProx} \\
                \midrule
                Local Epochs             & $5$             & $5$              \\
                Optimizer                & SGD             & SGD              \\
                Proximal Penalty ($\mu$) & --              & $0.01$           \\
                Status                   & Replicated      & Replicated       \\
                \bottomrule
            \end{tabular}
        \end{table}

        \column{0.52\textwidth}
        \centering
        \textbf{Validation Curves}
        \vspace{0.2cm}

        \begin{center}
            \setlength{\fboxrule}{0.5pt}
            \fbox{
                \begin{minipage}{0.9\textwidth}
                    \centering
                    \vspace{0.5cm}
                    \fontsize{7.5pt}{9.5pt}\selectfont
                    \texttt{[Placeholder: Loss/Accuracy Plots]} \\
                    \vspace{0.15cm}
                    \fontsize{6.5pt}{8.5pt}\selectfont
                    \textit{Source: Insert replicated loss/accuracy curves plotted from local logs}
                    \vspace{0.5cm}
                \end{minipage}
            }
        \end{center}
    \end{columns}
\end{frame}

% Slide 11: Standout Wrap-Up / Q&A
\begin{frame}[standout]
    Thank You!
    \vspace{0.5cm}

    \begin{center}
        \large Questions \& Answers
    \end{center}
\end{frame}

\end{document}
